\TemplateHeading
\templatefourrow{12.965}{课程名称}{计算机网络课程设计}{实验课时}{\experimenthours}
\templatefourrow{13.688}{实验项目}{IPv4子网划分与编址}{实验时间}{2026.11.17}
\templaterow{实验目的}{26}{
掌握IPv4地址、子网掩码与默认网关的配置方法，通过本地与远端Ping对照分析默认网关的作用；掌握VLSM子网划分、网络边界对齐和路由器接口编址方法，使用活动评分项与实际配置相互核对。
}
\templaterow{实验环境}{39}{
macOS；Cisco Packet Tracer 9.0.1。使用课程提供的实验6学生版四个活动文件，按各活动内地址表、主机数需求与指定接口进行配置。

指导教师：\teachername；课序号：06；班级全称：\classfullname；实验日期：\experimentdate。实验课时未核实，留空。
}
\templaterow{\experimentlabel}{96.977}{
\subheading{实验步骤与记录方法}
1. 题1沿用预设交换机、路由器及RIP配置，只配置三台PC的IPv4地址、掩码和网关；记录LAXpc不同网关条件下的完整Ping输出，包括失败结果，结束前恢复正确配置。

2. 题2配置三台PC和三台路由器的指定以太网接口，核对接口启用状态，并分别测试各PC到其默认网关的连通性。

3. 题3与题4先依据主机数从大到小划分九个子网，记录网络地址、掩码、广播地址、可用范围与十二个接口的实际编址，再完成指定接口配置。题4另核对HQ Serial0/2/0的时钟和启用状态。

4. 每题完成后运行Check Results，并展开Assessment Items，核对完成度与全部评分明细。分别记录配置截图、接口输出、Ping响应及统计，失败结果同样保留。

5. 报告按题1至题4组织实验结果、配置检查和问题分析。图片由本次真实界面截图裁剪，图下注明图号与图名；字段并排裁剪时明确标示省略的空白。课程提交一份报告，原始证据在本地留存。
}
